% -*- coding: utf-8 -*-
% XeLaTeX 编译
\documentclass[a3paper, landscape, 11pt, twocolumn]{article}

% 引入宏包
\usepackage{fontspec}
\usepackage[UTF8]{ctex}
\usepackage{geometry}
\usepackage{listings}
\usepackage{xcolor}
\usepackage{enumitem}
\usepackage{tabularx}
\usepackage{makecell}
\usepackage{fancyhdr}
\usepackage{amssymb}
\usepackage{lastpage}

% 设置页面边距 (A3横向双栏)
\geometry{left=2cm, right=2cm, top=2.5cm, bottom=2.5cm, columnsep=1.8cm}

% 颜色定义
\definecolor{codebg}{RGB}{248,248,248}
\definecolor{keywordcolor}{RGB}{0,0,150}
\definecolor{commentcolor}{RGB}{0,128,0}
\definecolor{stringcolor}{RGB}{163,21,21}

% Java代码块样式设置
\lstdefinestyle{JavaExam}{
    language=Java,
    basicstyle=\ttfamily\small,
    backgroundcolor=\color{codebg},
    keywordstyle=\color{keywordcolor}\bfseries,
    commentstyle=\color{commentcolor},
    stringstyle=\color{stringcolor},
    showstringspaces=false,
    breaklines=true,
    frame=single,
    rulecolor=\color{black!50},
    tabsize=4,
    xleftmargin=1em,
    xrightmargin=1em,
    aboveskip=1em,
    belowskip=1em
}

% 列表样式
\newlist{options}{enumerate}{1}
\setlist[options]{label=\Alph*., nosep, leftmargin=*, itemsep=0.5em}

% 页眉页脚设置
\pagestyle{fancy}
\fancyhf{}
\fancyhead[C]{\Large\bfseries 《Java程序设计》期末考试试卷}
\fancyfoot[C]{第 \thepage\ 页 共 \pageref{LastPage}\ 页}
\renewcommand{\headrulewidth}{0.4pt}

\begin{document}

% 试卷装订线与密封线模拟区
\twocolumn[{
			\begin{center}
				\vspace*{-1cm}
				{\huge\bfseries 2025—2026学年第一学期《Java程序设计》期末试卷}\\[1em]
				{\large 适用专业：计算机科学与技术、软件工程 \hspace{2em} 闭卷 \hspace{2em} 满分：150分 \hspace{2em} 考试时间：120分钟}\\[1em]
				\renewcommand{\arraystretch}{1.5}
				\begin{tabular}{|c|c|c|c|c|c|c|}
					\hline
					题号  & 一 (40分)        & 二 (20分)        & 三 (20分)        & 四 (30分)        & 五 (40分)        & 总分 (150分)      \\
					\hline
					得分  & \hspace{1.5cm} & \hspace{1.5cm} & \hspace{1.5cm} & \hspace{1.5cm} & \hspace{1.5cm} & \hspace{1.5cm} \\
					\hline
					阅卷人 &                &                &                &                &                &                \\
					\hline
				\end{tabular}\\[1.5em]
				{\large 姓名：\underline{\hspace{3cm}} \hspace{2em} 学号：\underline{\hspace{4cm}} \hspace{2em} 班级：\underline{\hspace{3cm}}}
				\vspace{1em}
				\hrule
				\vspace{1.5em}
			\end{center}
		}]

\section*{一、单项选择题（本大题共 20 小题，每小题 2 分，共 40 分）}
\textit{在每小题列出的四个备选项中只有一个是符合题目要求的，请将其代码填写在题后的括号内。}

\begin{enumerate}[label=\arabic*., itemsep=0.8em]
	\item 一个标准的Java源文件中，最多可以包含几个被 \texttt{public} 修饰的外部类？（ \quad ）
	      \begin{options}
		      \item 0 个
		      \item 1 个
		      \item 2 个
		      \item 不限制数量
	      \end{options}

	\item 执行数据类型转换代码 \texttt{byte b = (byte)130;} 后，变量 \texttt{b} 的值为（ \quad ）。
	      \begin{options}
		      \item 130
		      \item -126
		      \item 编译报错
		      \item 运行时异常
	      \end{options}

	\item 在Java位运算符中，关于 \texttt{>>} 和 \texttt{>>>} 的描述正确的是（ \quad ）。
	      \begin{options}
		      \item 两者完全没有区别
		      \item \texttt{>>} 是带符号右移，\texttt{>>>} 是无符号右移
		      \item \texttt{>>>} 是带符号右移，\texttt{>>} 是无符号右移
		      \item \texttt{>>} 用于整数，\texttt{>>>} 用于浮点数
	      \end{options}

	\item 若有一个二维数组 \texttt{int[][] arr = new int[3][5];}，获取其行数（即外层数组长度）的代码是（ \quad ）。
	      \begin{options}
		      \item \texttt{arr.length}
		      \item \texttt{arr[0].length}
		      \item \texttt{arr.size()}
		      \item \texttt{arr.length()}
	      \end{options}

	\item 下列访问修饰符中，访问权限最严格、封装级别最高的是（ \quad ）。
	      \begin{options}
		      \item \texttt{public}
		      \item \texttt{protected}
		      \item 默认 (缺省)
		      \item \texttt{private}
	      \end{options}

	\item 面向对象的方法重载（Overload）是指在一个类中（ \quad ）。
	      \begin{options}
		      \item 方法名相同，参数列表不同
		      \item 方法名相同，返回值类型不同
		      \item 参数列表相同，但修饰符不同
		      \item 子类覆盖父类的方法
	      \end{options}

	\item 多态机制下，当父类引用指向子类对象并调用一个被子类重写的方法时，实际执行的是（ \quad ）。
	      \begin{options}
		      \item 父类的方法
		      \item 子类重写后的方法
		      \item 编译报错
		      \item 抛出运行时异常
	      \end{options}

	\item 下列关于 \texttt{static} 关键字的描述中，\textbf{错误}的是（ \quad ）。
	      \begin{options}
		      \item 静态方法内部可以直接调用同一个类中的非静态方法
		      \item 静态方法可以直接通过“类名.方法名”的方式调用
		      \item 静态变量在内存中只有一份，被类的所有实例共享
		      \item 静态代码块在类被加载到内存时仅执行一次
	      \end{options}

	\item JDK 8 及以后的版本中，接口（Interface）内部允许存在带有具体实现的方法，这类方法需要使用哪种修饰符？（ \quad ）
	      \begin{options}
		      \item \texttt{default} 或 \texttt{static}
		      \item 只能是 \texttt{abstract}
		      \item \texttt{private}
		      \item \texttt{final}
	      \end{options}

	\item 局部内部类访问所在方法的局部变量时，该局部变量必须是（ \quad ）的（JDK 8以后可以隐式生效）。
	      \begin{options}
		      \item \texttt{static}
		      \item \texttt{final}
		      \item \texttt{public}
		      \item \texttt{abstract}
	      \end{options}

	\item 在方法声明处使用 \texttt{throws} 关键字的作用是（ \quad ）。
	      \begin{options}
		      \item 手动抛出一个具体的异常对象
		      \item 捕获并处理代码中的异常
		      \item 声明该方法执行过程中可能抛出的异常类型
		      \item 强制程序在此处停止运行
	      \end{options}

	\item 执行语句 \texttt{String s = "Java"; s = s + " SE";} 后，内存中发生了什么？（ \quad ）
	      \begin{options}
		      \item 原始的 "Java" 字符串对象的内容被修改了
		      \item 产生了一个全新的字符串对象 "Java SE"，变量s指向新对象
		      \item 发生编译错误，\texttt{String} 不支持 \texttt{+} 运算
		      \item 会抛出 \texttt{StringIndexOutOfBoundsException} 异常
	      \end{options}

	\item 自动装箱（Autoboxing）机制指的是（ \quad ）。
	      \begin{options}
		      \item 将基本数据类型自动转换为对应的包装类对象
		      \item 将包装类对象自动转换为对应的基本数据类型
		      \item 将字符串自动解析为数字
		      \item 将对象自动序列化为字节流
	      \end{options}

	\item 在Java集合框架中，若需要以“键值对”形式存储数据，并且希望遍历时保持插入的顺序，应该使用（ \quad ）。
	      \begin{options}
		      \item \texttt{HashMap}
		      \item \texttt{TreeMap}
		      \item \texttt{LinkedHashMap}
		      \item \texttt{Hashtable}
	      \end{options}

	\item 在泛型编程中，\texttt{List<? extends Number>} 代表的含义是（ \quad ）。
	      \begin{options}
		      \item 集合中的元素类型必须是 \texttt{Number} 或其子类
		      \item 集合中的元素类型必须是 \texttt{Number} 或其父类
		      \item 集合中的元素只能是 \texttt{Number} 类，不能是子类
		      \item 集合可以存放任何类型的对象
	      \end{options}

	\item 下列 I/O 流中，属于处理流（包装流）的是（ \quad ）。
	      \begin{options}
		      \item \texttt{FileInputStream}
		      \item \texttt{BufferedReader}
		      \item \texttt{FileOutputStream}
		      \item \texttt{FileReader}
	      \end{options}

	\item 在多线程编程中，调用某个线程对象的 \texttt{wait()} 方法时，当前线程必须已经持有（ \quad ）。
	      \begin{options}
		      \item 对象的锁（监视器）
		      \item I/O 资源的使用权
		      \item CPU的时间片
		      \item 数据库的连接句柄
	      \end{options}

	\item 图形界面编程中，使用事件适配器类（如 \texttt{WindowAdapter}）的主要目的是（ \quad ）。
	      \begin{options}
		      \item 提供比监听器接口更高的运行性能
		      \item 避免必须实现接口中的所有抽象方法，只需重写需要的方法
		      \item 让GUI界面自动支持多线程刷新
		      \item 让界面控件自动适应不同的屏幕分辨率
	      \end{options}

	\item JDBC操作中，正常关闭数据库资源的正确顺序是（ \quad ）。
	      \begin{options}
		      \item \texttt{Connection} $\rightarrow$ \texttt{Statement} $\rightarrow$ \texttt{ResultSet}
		      \item \texttt{ResultSet} $\rightarrow$ \texttt{Statement} $\rightarrow$ \texttt{Connection}
		      \item \texttt{Statement} $\rightarrow$ \texttt{ResultSet} $\rightarrow$ \texttt{Connection}
		      \item 关闭顺序完全无关紧要
	      \end{options}

	\item Java中要保证某个变量在多个线程间的可见性，但不保证复合操作的原子性，可以使用下列哪个关键字修饰？（ \quad ）
	      \begin{options}
		      \item \texttt{synchronized}
		      \item \texttt{transient}
		      \item \texttt{volatile}
		      \item \texttt{final}
	      \end{options}
\end{enumerate}

\section*{二、判断题（本大题共 10 小题，每小题 2 分，共 20 分）}
\textit{请判断下列说法的正误。正确的打“$\checkmark$”，错误的打“$\times$”。}

\begin{enumerate}[label=\arabic*., itemsep=0.8em]
	\item Java入口方法 \texttt{main} 的参数名必须叫 \texttt{args}，否则程序无法启动。 （ \quad ）
	\item 构造方法没有返回值类型，因此在构造方法体中绝对不能使用 \texttt{return;} 语句。 （ \quad ）
	\item Java中局部变量在使用之前必须进行显式的初始化，否则编译器会报错。 （ \quad ）
	\item 抽象类中必须至少包含一个抽象方法，否则无法被声明为抽象类。 （ \quad ）
	\item 两个对象的 \texttt{hashCode()} 返回值相同，则它们的 \texttt{equals()} 方法比较也必定返回 \texttt{true}。 （ \quad ）
	\item 受检异常（Checked Exception）在程序编写时必须通过 \texttt{try-catch} 捕获或通过 \texttt{throws} 声明抛出。 （ \quad ）
	\item \texttt{java.io.File} 类的对象可以直接调用方法向文件中写入字符串内容。 （ \quad ）
	\item \texttt{StringBuffer} 类的很多方法使用 \texttt{synchronized} 进行了修饰，因此它是线程安全的。 （ \quad ）
	\item 当线程调用 \texttt{Thread.sleep()} 方法进入睡眠时，该线程会释放它所持有的对象锁。 （ \quad ）
	\item \texttt{PreparedStatement} 不仅能够有效防止 SQL 注入攻击，并且在执行结构相同的SQL时效率更高。 （ \quad ）
\end{enumerate}

\section*{三、填空题（本大题共 5 小题，每空 2 分，共 20 分）}

\begin{enumerate}[label=\arabic*., itemsep=1em]
	\item Java程序的运行机制分为两个主要阶段，首先是 \underline{\hspace{4cm}} 阶段将源文件转为字节码，然后是 \underline{\hspace{4cm}} 阶段由JVM执行字节码。

	\item 在定义类时，如果不显式编写任何构造方法，编译器会自动提供一个 \underline{\hspace{3cm}} 的构造方法；如果类的成员变量被 \underline{\hspace{3cm}} 修饰，则它只能在当前类内部被访问。

	\item \texttt{super} 关键字除了能在子类中调用父类的构造方法，还可以用于访问父类被隐藏的 \underline{\hspace{3.5cm}} 和被子类重写的 \underline{\hspace{3.5cm}}。

	\item 在Java I/O包中，用于将字节输入流转换为字符输入流的转换流是 \underline{\hspace{4cm}}；为了提高字符流的读取效率，通常会外层包装一个 \underline{\hspace{4cm}} 流。

	\item 线程同步机制中，使用 \underline{\hspace{4cm}} 关键字可以对方法或代码块进行加锁控制；若要唤醒在对象监视器上等待的某一个线程，需调用该对象的 \underline{\hspace{4cm}} 方法。
\end{enumerate}

\section*{四、阅读程序，写出运行结果（本大题共 5 小题，每小题 6 分，共 30 分）}

\begin{enumerate}[label=\arabic*., itemsep=1em]
	\item \textbf{写出下列程序的输出结果：}
	      \begin{lstlisting}[style=JavaExam]
public class Demo1 {
    public static void main(String[] args) {
        int count = 0;
        outer: for (int i = 1; i <= 3; i++) {
            for (int j = 1; j <= 3; j++) {
                if (i * j > 4) break outer;
                count++;
            }
        }
        System.out.println("Count = " + count);
    }
}
\end{lstlisting}
	      \textbf{输出结果：}\underline{\hspace{8cm}}

	\item \textbf{写出下列程序的输出结果：}
	      \begin{lstlisting}[style=JavaExam]
class ClassA {
    public ClassA() { System.out.print("A"); }
}
class ClassB extends ClassA {
    public ClassB() { System.out.print("B"); }
    public ClassB(int x) { 
        this(); 
        System.out.print("C"); 
    }
}
public class Demo2 {
    public static void main(String[] args) {
        new ClassB(1);
    }
}
\end{lstlisting}
	      \textbf{输出结果：}\underline{\hspace{8cm}}

	\item \textbf{写出下列程序的输出结果：}
	      \begin{lstlisting}[style=JavaExam]
public class Demo3 {
    public static String getMessage() {
        try {
            int x = 10 / 0;
            return "TryBlock";
        } catch (Exception e) {
            return "CatchBlock";
        } finally {
            System.out.print("FinallyBlock ");
        }
    }
    public static void main(String[] args) {
        System.out.println(getMessage());
    }
}
\end{lstlisting}
	      \textbf{输出结果：}\underline{\hspace{8cm}}

	\item \textbf{写出下列程序的输出结果：}
	      \begin{lstlisting}[style=JavaExam]
public class Demo4 {
    public static void main(String[] args) {
        String s1 = "hello";
        String s2 = "hello";
        String s3 = new String("hello");
        System.out.println(s1 == s2);
        System.out.println(s1 == s3);
    }
}
\end{lstlisting}
	      \textbf{输出结果：}\underline{\hspace{8cm}}

	\item \textbf{写出下列程序的输出结果：}
	      \begin{lstlisting}[style=JavaExam]
class Point {
    int x;
    int y;
}
public class Demo5 {
    public static void modify(Point p) {
        p.x += 10;
        p = new Point();
        p.y = 20;
    }
    public static void main(String[] args) {
        Point pt = new Point();
        pt.x = 5;
        pt.y = 5;
        modify(pt);
        System.out.println(pt.x + ", " + pt.y);
    }
}
\end{lstlisting}
	      \textbf{输出结果：}\underline{\hspace{8cm}}
\end{enumerate}

\newpage
\section*{五、编程题（本大题共 2 小题，每小题 20 分，共 40 分）}

\textbf{1. 类的设计与封装（20分）}\\
请编写一个表示“图书（Book）”的类，具体要求如下：
\begin{itemize}
	\item (1) 包含三个私有属性：书名 (\texttt{String title})、作者 (\texttt{String author})、价格 (\texttt{double price})。
	\item (2) 提供一个包含三个参数的构造方法，用于初始化这三个属性。
	\item (3) 提供所有属性的 getter 和 setter 方法。特别要求在 \texttt{setPrice(double price)} 方法中进行数据校验：若传入的价格小于 0，则将价格强制设置为 0.0。
	\item (4) 提供一个 \texttt{public void display()} 方法，按格式打印图书的详细信息（如：“书名: Java编程, 作者: 张三, 价格: 55.5”）。
	\item (5) 编写一个测试类 \texttt{TestBook}，在 \texttt{main} 方法中创建一本书（书名《Java进阶》，作者"李四"，初始价格设定为 -20.0），然后调用 \texttt{display()} 方法查看输出效果。
\end{itemize}
\vspace{1em}
\textbf{【答题区】}
\vspace{12cm}

\textbf{2. 继承、抽象类与多态应用（20分）}\\
请编写一个“员工工资计算”的小型系统，要求如下：
\begin{itemize}
	\item (1) 定义一个抽象类 \texttt{Employee}（员工），包含私有属性：姓名 (\texttt{String name}) 和工号 (\texttt{String id})。提供带参构造方法、getter方法，并声明一个抽象方法 \texttt{public abstract double calculateSalary();} 用于计算工资。
	\item (2) 定义一个子类 \texttt{FullTimeEmployee}（全职员工）继承自 \texttt{Employee}。新增私有属性：基本月薪 (\texttt{double baseSalary})。提供相应的构造方法，并实现 \texttt{calculateSalary()} 方法（直接返回基本月薪）。
	\item (3) 定义另一个子类 \texttt{PartTimeEmployee}（兼职员工）继承自 \texttt{Employee}。新增私有属性：时薪 (\texttt{double hourlyRate}) 和工作小时数 (\texttt{int hours})。提供相应的构造方法，并实现 \texttt{calculateSalary()} 方法（工资 = 时薪 $\times$ 工作小时数）。
	\item (4) 编写测试类 \texttt{TestEmployee}，在 \texttt{main} 方法中利用多态创建一个全职员工对象和一个兼职员工对象。分别调用并打印这两名员工的姓名、工号及最终计算出的工资数值。
\end{itemize}
\vspace{1em}
\textbf{【答题区】}

% \newpage
% % ---------------------------------------------------------
% % 教师阅卷参考答案区（不需要打印给学生时可删除此页以下内容）
% % ---------------------------------------------------------
% \begin{center}
% 	{\huge\bfseries 教师阅卷参考答案及评分标准 (卷B)}
% \end{center}
% \vspace{1em}

% \textbf{一、单项选择题（每小题2分，共40分）}
% \begin{tabularx}{\textwidth}{XXXXX}
% 	1. B  & 2. B  & 3. B  & 4. A  & 5. D  \\
% 	6. A  & 7. B  & 8. A  & 9. A  & 10. B \\
% 	11. C & 12. B & 13. A & 14. C & 15. A \\
% 	16. B & 17. A & 18. B & 19. B & 20. C \\
% \end{tabularx}
% \textit{解析提示：}
% 第2题：byte范围是-128~127，130超范围，由于(byte)强转，截断后为-126。
% 第8题：静态方法内部不能直接调用非静态方法（无this对象）。
% 第17题：wait()必须在同步代码块/方法中调用，即必须持有对象锁。

% \vspace{1em}
% \textbf{二、判断题（每小题2分，共20分）}
% \begin{tabularx}{\textwidth}{XXXXX}
% 	1. $\times$ (参数名可以是任意合法标识符) & 2. $\times$ (可以用空的 return; 提前结束) & 3. $\checkmark$ & 4. $\times$ (抽象类可以没有抽象方法) & 5. $\times$ (哈希冲突时equals可能为false) \\
% 	6. $\checkmark$             & 7. $\times$ (File不能直接读写，需借助流)    & 8. $\checkmark$ & 9. $\times$ (sleep不会释放锁)  & 10. $\checkmark$                  \\
% \end{tabularx}

% \vspace{1em}
% \textbf{三、填空题（每空2分，共20分）}
% \begin{enumerate}[label=\arabic*., itemsep=0.5em]
% 	\item 编译 \quad 解释（或运行）
% 	\item 无参（或默认） \quad \texttt{private}
% 	\item 成员变量（或属性） \quad 成员方法（或方法）
% 	\item \texttt{InputStreamReader} \quad \texttt{BufferedReader}
% 	\item \texttt{synchronized} \quad \texttt{notify()}
% \end{enumerate}

% \vspace{1em}
% \textbf{四、阅读程序写结果（每小题6分，共30分）}
% \begin{enumerate}[label=\arabic*., itemsep=0.5em]
% 	\item \texttt{Count = 5} \quad \textit{(i=1跑3次，i=2时j=1跑1次,j=2跑1次,j=3时跳出大循环，共5次)}
% 	\item \texttt{ABC} \quad \textit{(先调父类无参构造打印A，再执行this()无参构造打印B，最后执行带参打印C)}
% 	\item \texttt{FinallyBlock CatchBlock} \quad \textit{(先执行finally输出打印，再将catch块的返回值返回输出)}
% 	\item \texttt{true}\\ \texttt{false} \quad \textit{(s1和s2指向常量池同一地址，s3是堆中新对象，错一行扣3分)}
% 	\item \texttt{15, 5} \quad \textit{(按引用传递，修改p.x有效，但形参p重新指向新对象后对原对象y无影响)}
% \end{enumerate}

% \vspace{1em}
% \textbf{五、编程题（每小题20分，共40分）}
% \\
% \textbf{第1题评分标准参考（20分）：}
% 属性定义及私有化(3分)；构造方法(3分)；getter和setter(4分)；setPrice逻辑校验(5分)；display方法及测试类(5分)。
% \begin{lstlisting}[style=JavaExam]
% class Book {
%     private String title;
%     private String author;
%     private double price;

%     public Book(String title, String author, double price) {
%         this.title = title;
%         this.author = author;
%         setPrice(price); // 借用set方法的校验逻辑
%     }
%     public String getTitle() { return title; }
%     public void setTitle(String title) { this.title = title; }
%     public String getAuthor() { return author; }
%     public void setAuthor(String author) { this.author = author; }
%     public double getPrice() { return price; }
    
%     public void setPrice(double price) {
%         if (price < 0) {
%             this.price = 0.0;
%         } else {
%             this.price = price;
%         }
%     }
%     public void display() {
%         System.out.println("书名: " + title + ", 作者: " + author + ", 价格: " + price);
%     }
% }
% public class TestBook {
%     public static void main(String[] args) {
%         Book book = new Book("Java进阶", "李四", -20.0);
%         book.display(); // 输出价格应为 0.0
%     }
% }
% \end{lstlisting}

% \vspace{1em}
% \textbf{第2题评分标准参考（20分）：}
% 抽象父类定义及抽象方法(5分)；FullTimeEmployee类定义及重写计算(5分)；PartTimeEmployee类定义及重写计算(5分)；测试类多态应用及输出(5分)。
% \begin{lstlisting}[style=JavaExam]
% abstract class Employee {
%     private String name;
%     private String id;
%     public Employee(String name, String id) {
%         this.name = name;
%         this.id = id;
%     }
%     public String getName() { return name; }
%     public String getId() { return id; }
%     public abstract double calculateSalary();
% }

% class FullTimeEmployee extends Employee {
%     private double baseSalary;
%     public FullTimeEmployee(String name, String id, double baseSalary) {
%         super(name, id);
%         this.baseSalary = baseSalary;
%     }
%     @Override
%     public double calculateSalary() { return baseSalary; }
% }

% class PartTimeEmployee extends Employee {
%     private double hourlyRate;
%     private int hours;
%     public PartTimeEmployee(String name, String id, double hourlyRate, int hours) {
%         super(name, id);
%         this.hourlyRate = hourlyRate;
%         this.hours = hours;
%     }
%     @Override
%     public double calculateSalary() { return hourlyRate * hours; }
% }

% public class TestEmployee {
%     public static void main(String[] args) {
%         Employee e1 = new FullTimeEmployee("王五", "F001", 8000.0);
%         Employee e2 = new PartTimeEmployee("赵六", "P001", 150.0, 40);
        
%         System.out.println(e1.getName() + "(" + e1.getId() + ") 工资: " + e1.calculateSalary());
%         System.out.println(e2.getName() + "(" + e2.getId() + ") 工资: " + e2.calculateSalary());
%     }
% }
% \end{lstlisting}

\end{document}